% =====================================================================
% VERSION ENSAE, REDUCTION SFE : conclusion.
% Phase 3 : copie de chapitres/14_conclusion.tex. Phase 4 (1er octobre 2026) : reecrite selon la
% regle n°2 de la mission et ramenee a deux pages, en suivant les consignes MS (situer le point
% d'arrivee, ce qui a pu etre realise, ce qui n'a pu l'etre, la suite possible). Rapport :
% reduction/chapitres/conclusion.md, qui liste ce qui est parti et ou.
% Reprise le 1er octobre 2026 tard le soir (« une ecriture plus humaine ») : registre cale sur
% des memoires d'actuariat publies avant 2022 (style/registre_humain.md), longueur inchangee, la
% derniere page etant pleine. Fond inchange.
% Lu par main_ensae.tex seul : ne jamais le faire lire par main_institut.tex.
% =====================================================================

\chapter*{Conclusion}
\phantomsection
\addcontentsline{toc}{chapter}{Conclusion}
\markboth{Conclusion}{}

\section*{Résultats}
Les montants concernent le risque opérationnel d'origine informatique du périmètre financier étudié, et non celui d'un assureur.
Le besoin de capital passe de $6\,049$ à $20\,188$~M\euro{} entre les états conforme et non conforme des quatre canaux (fréquence, détection, propagation, accumulation),
soit un facteur $3{,}34$. L'écart de $14\,139$~M\euro{} compare deux états posés, sans mesurer un effet observé de DORA.
Son écart-type est de $734$~M\euro{} entre quatre graines de simulation.

Les canaux isolés en expliquent
$9\,138$~M\euro{} et leur interaction $5\,001$~M\euro{}, soit $35\,\%$ de l'écart, ce que donnerait à
peu près un produit des facteurs. Un plan de remédiation doit donc citer l'effet d'un canal fermé
depuis l'état non conforme. Relâcher la seule fréquence fait passer le capital de l'état conforme à
$10\,377$~M\euro{}, et toute cible au-delà de $11\,413$~M\euro{} exige ce canal.

La cascade distingue des ordres d'atteinte qu'un score statique des seuls piliers touchés ne représente pas :
$22$ des $26$ ensembles de piliers changent de criticité (note de $1$ à $10$) selon l'ordre d'atteinte.
La normalisation de Leontief garantit la stabilité de la contagion, cycles compris. Un processus de Hawkes
ne se distingue pas de la cascade à la résolution disponible, et nous retenons la cascade par parcimonie.

La queue de sévérité ($\hat\xi = 0{,}5954$ au seuil publié) et la fréquence d'entrée sont calibrées sur
les données ; la propagation et l'accumulation sont posées. Hors échantillon, la binomiale négative
couvre $91{,}7\,\%$ des années contre $75\,\%$ pour le Poisson, et la forme de la queue n'est pas rejetée
($p = 0{,}4087$), mais douze années donnent peu de puissance. Le niveau de la charge annuelle est rejeté
(Kolmogorov-Smirnov, $p = 0{,}0042$), en raison d'une dérive mesurée de l'échelle de sévérité. La
Formule Standard n'exprime pas cet effet.

% SOURCES-SCRIPTS: 28 39 43 68 92 47 89 91 85 100 74 32

\section*{Non-identifiabilité de la direction}
La direction de la contagion n'est pas identifiable sur les données publiques, et les trois sources
étudiées (chronologie américaine, SAS OpRisk, VERIS) échouent pour des raisons distinctes. La meilleure expérience naturelle disponible,
le choc de chaîne d'approvisionnement MOVEit, a un effet net nul une fois un placebo retranché. Le
schéma VERIS, standard de description d'incidents de la profession, prévoit bien les deux champs
manquants, mais leur intersection n'est renseignée que dans \emph{un} incident sur $10\,591$.

Le socle sans
contagion ne dépend pas de la direction, seule grandeur non identifiable ; il garde des paramètres posés
(surdispersion, accumulation), voir le chapitre~\ref{ch:robustesse}. Au-dessus de lui, dans le calage de référence (sévérité OpRisk), le capital du périmètre financier étudié se situe entre
$6\,858$ et $8\,697$~M\euro{} sur les $1\,024$ sommets de l'ensemble admissible. Cette plage, large de $1\,839$~M\euro{}, est calculée aux sommets et n'intègre pas l'incertitude sur la partie symétrique posée par l'expert.

% SOURCES-SCRIPTS: 30 34 28 110

\section*{Conséquences pour le régulateur et pour l'entité}
Ces trois échecs indiquent la donnée qui rendrait la contagion mesurable : un registre d'incidents.
Il horodaterait la survenance à la semaine et classerait chaque incident selon le domaine de
contrôle en cause, plutôt que selon le vecteur d'attaque. Il consoliderait une cause commune en un
événement unique et consignerait l'ensemble des piliers touchés. La loi jointe de ces ensembles
suffit à distinguer une matrice de sa transposée.

Les deux premiers champs existent déjà dans le schéma VERIS, et la notification d'incident majeur
prévue par DORA peut en rendre le remplissage obligatoire, ce qu'une base déclarative volontaire ne
peut pas. Sur une grille commune, fixer le sens des six dépendances documentées par sept rapports post-incident réduit la plage de $32{,}7\,\%$ en moyenne,
de $22$ à $44\,\%$ selon la graine. Ce résultat dépend de la grille et du barème d'expert. Les dépendances les plus utiles à documenter, dont celle entre
gouvernance et gestion des incidents, en lèveraient chacune de $11$ à $13\,\%$ en moyenne, sans que leur ordre soit
résolu.

% SOURCES-SCRIPTS: 66 55 117

Remédier, détenir ou transférer le risque se comparent ensuite en euros par an, avec une prime de réassurance
calculée sur un ratio de sinistralité de $48{,}5\,\%$. Pour l'entité de $15\,000$~M\euro{} de provisions,
dont le besoin de capital ($169$~M\euro{}) est une extrapolation d'ordre de grandeur, le coût annuel
de détention représente plus de deux fois la sinistralité attendue. Avec le traité en excès de perte idéalisé,
le coût de financement baisse d'environ deux tiers, et de plus de moitié pertes conservées comprises.
La capacité du marché, les exclusions et le risque de contrepartie réduisent ce gain.



% SOURCES-SCRIPTS: 58 65 67

\section*{Limites}
L'élicitation structurée a été préparée et diffusée, mais elle est restée sans réponse (supplément,
\S\ref{ch:elicitation}). La lecture centrale de la direction repose donc sur un barème d'expert
qualitatif. Deux hypothèses pèsent le plus sans être testables sur la donnée : les niveaux de
propagation, posés, et l'additivité des coûts entre piliers touchés. L'élasticité de cette dernière atteint $0{,}95$ à
l'état non conforme contre $0{,}39$ à l'état conforme (\S\ref{sec:additivite-couts}).
Dans les deux cas, le signe et l'ordre de l'écart tiennent sur toute la plage étudiée, mais pas
son amplitude.

Sur les sensibilités étudiées, le signe du surcoût et les deux piliers les plus coûteux, la gouvernance et la gestion des
prestataires tiers, sont conservés. L'ordre de remédiation entre ces deux piliers dépend en revanche du
niveau d'émission de la gouvernance : il tient tant que celle-ci n'est dégradée que du
quart, et bascule à la moitié. Le niveau absolu du capital, du secteur comme de l'entité, est illustratif : l'incertitude de queue impose de le présenter en distribution.

% SOURCES-SCRIPTS: 74 64 65

\section*{Ouvertures}
L'élicitation structurée, si un panel y répondait, resserrerait la bande de capital. Une mesure de la maîtrise des risques par pilier permettrait de remplacer le gain de propagation
posé par une grandeur observée.
